%%%%%%%%%%%%%%%%%%%%%%%%%%%%%%%%%%%%%%%%%%%%%%%%%%%%%%%%%%%%%%%
% 
% A transform involving Chebyshev polynomials and its inversion formula,
% by
% \'Oscar Ciaurri, Luis Navas and Juan L. Varona
% 
% COMIENZO ESCRITURA: 11-octubre-2004
% ULTIMA MODIFICACION: 5-septiembre-2005
%%%%%%%%%%%%%%%%%%%%%%%%%%%%%%%%%%%%%%%%%%%%%%%%%%%%%%%%%%%%%%%

\documentclass[10pt,a4paper]{amsart}

%\setlength{\textheight}{43pc} 
%\setlength{\textwidth}{28pc}

%\renewcommand{\baselinestretch}{1.8} %% 2.0 is double spacing
%\renewcommand{\baselinestretch}{1.74} %% 2.0 is double spacing

%\usepackage[notcite,notref]{showkeys}

% \newcommand{\datename}{\textit{Date}:} %% <- Original
\renewcommand{\datename}{THIS PAPER HAS BEEN PUBLISHED IN:}
\date{\textit{J.\@ Math.\@ Anal.\@ Appl.\@} \textbf{323} (2006), 57--62.}

\theoremstyle{plain}
\newtheorem{theorem}{Theorem}
\newtheorem{proposition}{Proposition}

\theoremstyle{remark}
\newtheorem*{remark}{Remark} % without number

\renewcommand{\Re}{\operatorname{Re}}

\newcommand{\N}{\mathbb{N}}
\newcommand{\Q}{\mathbb{Q}}
\newcommand{\R}{\mathbb{R}}
\newcommand{\C}{\mathbb{C}}
\newcommand{\Z}{\mathbb{Z}}
\newcommand{\numdiv}{d} % number of divisors. % Alternative: \tau
% But \tau is the Ramanujan function.

%%%%%%%%%%%%%%%%%%%%%%%%%%%%%%%%%%%%%%%%%%%%%%%%%%%%%%
\begin{document}
%%%%%%%%%%%%%%%%%%%%%%%%%%%%%%%%%%%%%%%%%%%%%%%%%%%%%%


\title[A transform involving Chebyshev polynomials]{A transform involving
Chebyshev polynomials and its inversion formula}

% \author{\'Oscar Ciaurri, Luis Navas and Juan L. Varona}
% \email{oscar.ciaurri@dmc.unirioja.es, navas@usal.es,
% jvarona@dmc.unirioja.es}

\author{\'Oscar Ciaurri}
\address{
        Departamento de Matem\'aticas y Computaci\'on \\
        Universidad de La Rioja                       \\
        Calle Luis de Ulloa s/n                       \\
        26004 Logro\~no, Spain}
\email{oscar.ciaurri@dmc.unirioja.es}

\author{Luis M. Navas}
\address{
        Departamento de Matem\'aticas                \\
        Universidad de Salamanca                     \\
        Plaza de la Merced 1-4                       \\
        37008 Salamanca, Spain}
\email{navas@usal.es}

\author{Juan L. Varona}
\address{
        Departamento de Matem\'aticas y Computaci\'on \\
        Universidad de La Rioja                       \\
        Calle Luis de Ulloa s/n                       \\
        26004 Logro\~no, Spain}
\email{jvarona@dmc.unirioja.es}
\urladdr{http://www.unirioja.es/dptos/dmc/jvarona/welcome.html}

\thanks{Research of the first and third authors supported by grant
BFM2003-06335-C03-03 of the~DGI}

\keywords{Chebyshev polynomials, M\"obius function, M\"obius transform,
Dirichlet convolution, inversion formula}

\subjclass[2000]{Primary 11A25; Secondary 33C45}

%\date{February 8, 2005}

%%%%%%%%%%%%%%%%%%%%%%%%%%%%%%%%%%%%%%%%%%%%%%%%%%%%%%%%%%%%%%%%%%%%%%%%%%%%%
\begin{abstract}
We define a functional analytic transform involving the Chebyshev polynomials
$T_{n}(x),$ with an inversion formula in which the M\"obius function $\mu(n)$
appears. If $s \in \mathbb{C}$ with $\operatorname{Re}(s) > 1,$ then given a
bounded function from $[-1,1]$ into $\mathbb{C}$, or from $\mathbb{C}$ into
itself, the following inversion formula holds:
$$
  g(x) = \sum_{n=1}^\infty \frac{1}{n^s}\, f(T_n(x))
$$
if and only if
$$
  f(x) = \sum_{n=1}^\infty \frac{\mu(n)}{n^s}\, g(T_n(x)).
$$
Some other similar results are given.
\end{abstract}
%%%%%%%%%%%%%%%%%%%%%%%%%%%%%%%%%%%%%%%%%%%%%%%%%%%%%%%%%%%%%%%%%%%%%%%%%%%%%


%%%%%%%%%%
\maketitle
%%%%%%%%%%


%%%%%%%%%%%%%%%%%%%%%%%%%%%%%%%%%%%%%%%%%%%%%%%%%%%%%%%%%%%%%%%%%%%%%%%%%%%%%
\section{Introduction and main results}
%%%%%%%%%%%%%%%%%%%%%%%%%%%%%%%%%%%%%%%%%%%%%%%%%%%%%%%%%%%%%%%%%%%%%%%%%%%%%

If we have an arithmetical function $\alpha: \N \to \R$ (or $\C$) and a function
$f: (0,\infty) \to \R$ (or $\C$), we can define a new function
$g = \alpha  \circ f$ by taking
\begin{equation}
\label{eq:convocirc}
  g(x) = (\alpha \circ f)(x) = \sum_{n=1}^\infty \alpha(n) f\left(\frac xn\right),
  \qquad x \in (0,\infty).
\end{equation}
Moreover, let us suppose that $\alpha$ is invertible with respect to Dirichlet
convolution (this happens if and only if $\alpha(1) \ne 0$).
Then, it is well known that we have the inversion formula
$f = \alpha^{-1} \circ g$.
A typical case is when $\alpha$ is a completely multiplicative function;
in this case, $\alpha^{-1}(n) = \mu(n)\alpha(n)$,
where $\mu(n)$ is the M\"obius function.
Thus, we have
\begin{equation*}
  f(x) = \sum_{n=1}^\infty \mu(n)\alpha(n) g\left(\frac xn\right),
  \qquad x \in (0,\infty)
\end{equation*}
(see, for instance, \cite{Ha-Wr}).
A common example of a completely multiplicative function is
$\alpha(n) = n^{-s}$, $s \in \C,$ which gives rise to Dirichlet series.

In this paper we present a new transform/inverse pair in which
both the Chebyshev polynomials $\{T_n(x)\}_{n=1}^\infty$
and the M\"obius function $\mu(n)$ appear.
The Chebyshev polynomials satisfy many identities and orthogonal conditions,
but for our purposes only the property
\begin{equation}
\label{eq:TnTm=Tnm}
  T_m(T_n(x)) = T_{mn}(x)
\end{equation}
is essential.
For $x \in [-1,1]$, this formula is clear from $T_k(x) = \cos(k\arccos x)$ and,
for $x \in \C$, it follows by analytic continuation.
% It is interesting to note that, if a family of polynomials satisfy an identity
% similar to~\eqref{eq:TnTm=Tnm}, upon a linear change of variable
% this family must be $\{z^n\}$ or the Chebyshev polynomials;
% see~\cite[]{Ri} for details.
It is interesting to note that, up to a linear change of variable,
$\{x^n\}$ and the Chebyshev polynomials are the unique families of
polynomials that satisfy an identity similar to~\eqref{eq:TnTm=Tnm}
(see~\cite[Chapter~4, Th.~4.4]{Ri} for details);
in particular, similar inversion formulas can be given for
expansions in $\{x^n\}$ on $[0,1]$.
%	There is a huge quantity of mathematical literature related to
%	Chebyshev polynomials; however, it seems that the kind of transform
%	and its inverse that we show have not been noticed before.

Prior to continuing, let us establish our notation. For $x \in [-1,1]$,
we have $T_n(x) = \cos(n\arccos x)$ so $T_n: [-1,1] \to [-1,1]$.
But we can also consider $T_n$ both as $T_n: \R \to \R$ or $T_n: \C \to \C.$
Moreover, $T_n(x) \in \Z[x]$, so we also have
$T_n: \Z \to \Z$ and $T_n: \Q \to \Q$.
Let us then use $\Delta$ to denote $[-1,1]$, $\R$, $\C$, $\Z$, or~$\Q$, accordingly.
Thus, for functions $f$ of type $f: \Delta \to \R$ (or~$\C$),
the composition $f(T_n(x))$ is well defined for every~$n$.
(Some perhaps more ``esoteric'' choices can be taken into account for $\Delta$,
such us $[1,\infty)$, $\N$, or the algebraic numbers.)

The main result of this paper is the following:

\begin{theorem}
\label{th:f<->gChe}
Let $s \in \C$ with $\Re(s)>1$
and $\Delta = [-1,1]$, $\R$, $\C$, $\Z$, or~$\Q$.
If $f$ is a bounded function defined on $\Delta$, then the series
\begin{equation}
\label{eq:f->gChe}
  g(x) = \sum_{n=1}^\infty \frac{1}{n^s}\, f(T_n(x)), \qquad x \in \Delta,
\end{equation}
is absolutely convergent, the function $g$ is bounded, and we can recover $f$ as
\begin{equation}
\label{eq:g->fChe}
  f(x) = \sum_{n=1}^\infty \frac{\mu(n)}{n^s}\, g(T_n(x)), \qquad x \in \Delta.
\end{equation}
Conversely, if we have a bounded function $g$ on $\Delta$, the function
$f$ defined as in~\eqref{eq:g->fChe} is bounded and fulfills~\eqref{eq:f->gChe}.
\end{theorem}

%	\begin{remark}
%	As operators on the Banach space $L^{\infty}(\Delta)$, in the situation of
%	this theorem one has the estimates
%	\begin{align*}
%	  \|\mathcal{T}_sf\|_\infty      &\le \zeta(\Re(s)) \|f\|_\infty, \\
%	  \|\mathcal{T}_s^{-1}g\|_\infty &\le \zeta(\Re(s)) \|g\|_\infty,
%	\end{align*}
%	where $\zeta$ is the Riemann zeta function.
%	\end{remark}

We also study some other conditions that yield similar results.
In particular, in Section~\ref{sec:general} we give a more general approach to our inversion formula.

%%%%%%%%%%%%%%%%%%%%%%%%%%%%%%%%%%%%%%%%%%%%%%%%%%%%%%%%%%%%%%%
\section{The transform and the inversion formula: proof of the main theorem}
\label{sec:transform}
%%%%%%%%%%%%%%%%%%%%%%%%%%%%%%%%%%%%%%%%%%%%%%%%%%%%%%%%%%%%%%%

Let us begin by defining an operation $\odot$ similar to the~$\circ$
in~\eqref{eq:convocirc}, but properly adapted to our circumstances.
Given a function $f$ on $\Delta$ and an arithmetical function
$\alpha: \N \to \R$ (or $\C$), we define the transform
\begin{equation}
\label{eq:convoodot}
  g(x) = (\alpha \odot f)(x) = \sum_{n=1}^\infty \alpha(n) f(T_n(x)),
\end{equation}
provided that the series converges.

Let us suppose that we have another arithmetical function $\beta: \N \to \R$ (or $\C$)
that is inverse to $\alpha$ with respect to Dirichlet convolution,
i.e., $\alpha * \beta = \delta$ with $\delta(1) = 1$ and $\delta(n)=0$ for $n>1$.
% (usually, we denote $\beta = \alpha^{-1}$)
Let us calculate $(\beta \odot g)(x)$, at least formally, from~\eqref{eq:convoodot}. If the formal manipulations that follow are analytically justified, we can reorder series, group the terms such than $nm=k$, use \eqref{eq:TnTm=Tnm}, $\alpha * \beta = \delta$, and $T_1(x) = x$, so
\begin{equation}
\label{eq:inv-formal}
\begin{aligned}
  (\beta \odot g)(x)
  &= \sum_{n \in \N} \beta(n) (\alpha \odot f)(T_n(x)) \\
  &= \sum_{n \in \N} \beta(n) \sum_{m \in \N} \alpha(m) f(T_m(T_n(x))) \\
  &= \sum_{n,m \in \N} \beta(n) \alpha(m) f(T_{mn}(x)) \\
  &= \sum_{k \in \N} \left(\sum_{nm = k} \beta(n) \alpha(m) \right)
    f(T_k(x)) \\
  &= \sum_{k \in \N} (\alpha*\beta)(k) f(T_k(x)) \\
%  &= \sum_{k \in \N} \delta(k) f(T_k(x)) \\
  &= f(x).
\end{aligned}
\end{equation}
Thus, we have found the inversion formula.
It remains to determine conditions under which the series that define $(\alpha \odot f)(x)$ and $(\beta \odot g)(x)$ converge and the manipulations in~\eqref{eq:inv-formal} can be justified.

Some simple assumptions guaranteeing this are the following:

\begin{proposition}
\label{prop:bounded}
Let $\alpha$ and $\beta$ be two arithmetical functions related by $\alpha*\beta = \delta$, and such that $\sum_{n=1}^\infty |\alpha(n)| < \infty$ and $\sum_{n=1}^\infty |\beta(n)| < \infty$; let $\Delta$ be $[-1,1]$, $\R$, $\C$, $\Z$, or~$\Q$.
If $f$ is a bounded function defined on $\Delta$, then the series
\begin{equation}
\label{eq:f->ga}
  g(x) = \sum_{n=1}^\infty \alpha(n) f(T_n(x)), \qquad x \in \Delta,
\end{equation}
is absolutely convergent, the function $g$ is bounded by
\begin{equation}
\label{eq:boundg}
  \sup_{x \in \Delta} |g(x)|
  \le \Big( \sum_{n=1}^\infty |\alpha(n)| \Big) \sup_{x \in \Delta} |f(x)|,
\end{equation}
and we can recover $f$ as
\begin{equation}
\label{eq:g->fb}
  f(x) = \sum_{n=1}^\infty \beta(n) g(T_n(x)), \qquad x \in \Delta.
\end{equation}
Conversely, if we have a bounded function $g$ on $\Delta$, the function
$f$ defined as in~\eqref{eq:g->fb} is bounded in a similar way and
fulfills~\eqref{eq:f->ga}.
\end{proposition}

%	\begin{proof}
%	It is clear that, under these conditions, the series
%	$g(x) = (\alpha \odot f)(x)$ converges absolutely, and $g$ satisfies~\eqref{eq:boundg}.
%	Then, also $(\beta \odot g)(x)$ converges absolutely. The same happens with the series 
%	involved in~\eqref{eq:inv-formal}, so the rearrangement of the terms is justified.

%	Of course, we can also start from $g$, define $f = \beta \odot g$ and recover
%	$g = \alpha \odot g$ in a similar way.
%	\end{proof}

With this, we have,

\begin{proof}[Proof of Theorem~\ref{th:f<->gChe}]
In the proposition, take $\alpha(n) = \alpha_s(n) = n^{-s}$,
which is a completely multiplicative function whose inverse is
$\alpha^{-1}(n) = \mu(n) n^{-s}$.
As $\Re(s)>1$, it follows that $\sum_{n=1}^\infty |a(n)|
= \sum_{n=1}^\infty |n^{-s}| = \zeta(\Re(s))$, where
$\zeta(s)$ denote the Riemann's zeta function.
The inversion part is similar.
\end{proof}

%%%%%%%%%%%%%%%%%%%%%%%%%%%%%%%%%%%%%%%%%%%%%%%%%%%%%%%%%%%%%%%
\subsection*{Another example}
%%%%%%%%%%%%%%%%%%%%%%%%%%%%%%%%%%%%%%%%%%%%%%%%%%%%%%%%%%%%%%%

Let us consider the Liouville function $\lambda(n),$ defined by
\begin{equation*}
%\label{eq:lambda}
  \lambda(n) =
  \begin{cases}
    1, & \mbox{if } n=1,\\
    (-1)^{a_1+\cdots+a_k}, & \mbox{if } n = p_1^{a_1}\cdots p_k^{a_k}
  \end{cases}
\end{equation*}
(where $p_1^{a_1}\cdots p_k^{a_k}$ denotes the decomposition of~$n$
into prime factors). $\lambda(n)$ is a completely multiplicative whose
inverse function is $\lambda^{-1}(n) = \mu(n)\lambda(n) = |\mu(n)|$.
Then, in a similar way to Theorem~\ref{th:f<->gChe}, for $\Re(s)>1$ and bounded
functions, we have
\begin{equation*}
%\label{eq:f->gChe-lambda}
  g(x) = \sum_{n=1}^\infty \frac{\lambda(n)}{n^s} f(T_n(x)), \qquad x \in \Delta,
\end{equation*}
if and only if
\begin{equation*}
%\label{eq:g->fChe-lambda}
  f(x) = \sum_{n=1}^\infty \frac{|\mu(n)|}{n^s} \, g(T_n(x)), \qquad x \in \Delta.
\end{equation*}


%%%%%%%%%%%%%%%%%%%%%%%%%%%%%%%%%%%%%%%%%%%%%%%%%%%%%%%%%%%%%%%
\section{A more general approach}
\label{sec:general}
%%%%%%%%%%%%%%%%%%%%%%%%%%%%%%%%%%%%%%%%%%%%%%%%%%%%%%%%%%%%%%%

The assumptions in Proposition~\ref{prop:bounded} are very demanding.
Here we study other general conditions under which the transformation
formula holds.

For an arithmetical function~$\rho$, we say that $f \in L(\Delta,\rho)$ if
\begin{equation*}
  \sum_{n=1}^\infty |\rho(n) f(T_n(x))| < \infty
  \quad \forall x \in \Delta
\end{equation*}
(recall that we are using $\Delta$ to denote $[-1,1]$, $\R$, $\C$, $\Z$ or~$\Q$).
In particular, $f \in L(\Delta,\alpha)$ means that~\eqref{eq:convoodot}
converges absolutely for every $x \in \Delta$.

% \delta ya la he definido antes, pero por recordar...
Once again we use the arithmetical function $\delta$ defined by
$\delta(1) = 1$ and $\delta(n) = 0$ for all $n > 1$.
The relation $\delta \odot f = f$ follows easily from $T_1(x) = x$.
%indeed, for every $x \in \Delta$, we have
%\begin{equation*}
%  (\delta \odot f) (x) = \sum_{n=1}^\infty \delta(n) f(T_n(x))
%  = 1 \cdot f(x) + 0 = f(x).
%\end{equation*}

Analogously to the mixed associative property between $\circ$
and Dirichlet convolution~$*$, we have the following version between $\odot$ and~$*$.
% (note that $\alpha * \beta$ is the ordinary
% Dirichlet convolution, not a variant related with~$\odot$)
The proof is straightforward, because the absolute convergence
allows the rearrangement of the sums.

\begin{proposition}
\label{pro:mixed}
Let $\alpha, \beta$ be two arithmetical functions, $f: \Delta \to \R$
(or~$\C$), and suppose at a given $x \in \Delta$
\begin{equation}
\label{eq:conabs}
  \sum_{n,m \in \N} |\alpha(n) \beta(m) f(T_{nm}(x))|
  = \sum_{k \in \N} (|\alpha| * |\beta|)(k) |f(T_k(x))|
  < \infty.
\end{equation}
Then, all the series involved in the definitions of 
$(\alpha \odot (\beta \odot f))(x)$ and $((\alpha * \beta) \odot f)(x)$
are absolutely convergent, and
\begin{equation*}
  (\alpha \odot (\beta \odot f))(x)
  = ((\alpha * \beta) \odot f)(x).
\end{equation*}
In particular, if $f \in L(\Delta,|\alpha|*|\beta|)$, then
$\alpha \odot (\beta \odot f) = (\alpha * \beta) \odot f$.
\end{proposition}

%	\begin{proof}
%	By applying property~\eqref{eq:TnTm=Tnm}, and grouping together
%	the $n,m\in\N$ such that $nm=k$, we have
%	\begin{align*}
%	  (\alpha \odot (\beta \odot f))(x)
%	  &= \sum_{n \in \N} \alpha(n) (\beta \odot f)(T_n(x)) \\
%	  &= \sum_{n \in \N} \alpha(n) \sum_{m \in \N} \beta(m) f(T_m(T_n(x))) \\
%	  &= \sum_{n,m \in \N} \alpha(n) \beta(m) f(T_{nm}(x)) \\
%	  &= \sum_{k \in \N} \left(\sum_{nm = k} \alpha(n) \beta(m) \right)
%	    f(T_k(x)) \\
%	  &= \sum_{k \in \N} (\alpha*\beta)(k) f(T_k(x))
%	  = ((\alpha * \beta) \odot f)(x),
%	\end{align*}
%	where the previous series are absolutely convergent by~\eqref{eq:conabs},
%	and thus the rearrangement of the sums is justified.
%	\end{proof}

In this general context, the inversion formula becomes

\begin{proposition}
\label{pro:f<->g-alpha}
Let $\alpha$ be an arithmetical funcion with Dirichlet convolution
inverse~$\alpha^{-1}$. Given a function $f: \Delta \to \R$ (or $\C$),
with $f \in L(\Delta,|\alpha|*|\alpha^{-1}|)$, the transform
$g(x) = (\alpha \odot f)(x)$
is defined for all $x \in \Delta$. Moreover, if 
$g \in L(\Delta,|\alpha|*|\alpha^{-1}|)$, then
$f(x) = (\alpha^{-1} \odot g)(x)$
for all $x \in \Delta$.
\end{proposition}

\begin{proof}
By Proposition~\ref{pro:mixed},
\begin{equation*}
  \alpha^{-1} \odot g = \alpha^{-1} \odot (\alpha \odot f)
  = (\alpha^{-1} * \alpha) \odot f = \delta \odot f = f.
\end{equation*}
For the second part, recall that
$|\alpha|*|\alpha^{-1}| = |\alpha^{-1}|*|\alpha|$.
\end{proof}

% AN IMMEDIATE COMMENT:
%	It is also possible to deduce $g = \alpha \odot f$ starting from
%	$f = \alpha^{-1} \odot g$ simply by replacing $f$ with $g$ and
%	$\alpha$ by $\alpha^{-1}$ in Proposition~\ref{pro:f<->g-alpha}.

In general, it does not seem easy to check that the condition
$f \in L(\Delta,|\alpha|*|\alpha^{-1}|)$ implies
$g \in L(\Delta,|\alpha|*|\alpha^{-1}|)$; this---if true---would mean
that the inversion formula $\alpha^{-1} \odot g$ is defined without
this extra hypothesis.
%	Of course the same phenomenon occurs with the ordinary Fourier transform:
%	for every $f \in L^{1}(\R)$, the transform $\widehat{f}$ is defined,
%	but it is not always true that $\widehat{f} \in L^{1}(\R)$,
%	so the inverse transform could be applied.

The following special case of Proposition~\ref{pro:f<->g-alpha} has special interest:

\begin{proposition}
\label{pro:f<->g-acm}
Let $\alpha$ be a completely multiplicative arithmetical function,
$f: \Delta \to \R$ (or $\C$), and suppose that
$f \in L(\Delta,\alpha\numdiv)$ (where $\numdiv(n)$ is the number
of divisors of~$n$).
Then
\begin{equation*}
  g(x) = \sum_{n \in \N} \alpha(n) f(T_n(x))
\end{equation*}
is defined for all $x \in \Delta$. Moreover, if $g \in L(\Delta,\alpha\numdiv)$, then
\begin{equation*}
  f(x) = \sum_{n \in \N} \mu(n) \alpha(n) g(T_n(x))
\end{equation*}
for all $x \in \Delta$.
\end{proposition}

\begin{proof}
If $\alpha$ is completely multiplicative, then $\alpha^{-1}(n) =
\mu(n)\alpha(n)$. Moreover,
\begin{align*}
  (|\alpha|*|\alpha^{-1}|)(k) &= \sum_{nm=k} |\alpha(n)\mu(m)\alpha(m)| \\
  &\le \sum_{nm=k} |\alpha(nm)| = \numdiv(k)|\alpha|(k) = |\numdiv(k)\alpha(k)|,
\end{align*}
so the hypothesis $f \in L(\Delta,\alpha\numdiv)$ allows us to
apply Proposition~\ref{pro:f<->g-alpha}. The same holds with respect
to $g \in L(\Delta,\alpha\numdiv)$.
\end{proof}

\begin{remark}
As commented previously, it does not seem easy to check if
$g \in L(\Delta,\alpha\numdiv)$ given that
$f \in L(\Delta,\alpha\numdiv)$. However, we claim that something weaker is true:
\begin{equation*}
  f \in L(\Delta,\alpha\numdiv^2) \;\Rightarrow\;
  g \in L(\Delta,\alpha\numdiv).
\end{equation*}
To prove this, take into account that $\numdiv(n) \le \numdiv(k)$ when $n \mathrel{|} k$,
and notice also that $|\alpha|$ is completely multiplicative. 
Thus
\begin{align*}
\sum_{n \in \N} \numdiv(n)|\alpha(n) g(T_n(x))| 
&= \sum_{n \in \N} \numdiv(n)|\alpha(n)|
  \left|\sum_{m \in \N} \alpha(m) f(T_m(T_n(x))) \right| \\
&\le \sum_{k \in \N} \numdiv(k)
  \left(\sum_{nm = k} |\alpha(n) \alpha(m)| \right) |f(T_k(x))| \\
&= \sum_{k \in \N} \numdiv(k)^2 |\alpha(k)| |f(T_k(x))| < \infty
\end{align*}
since $f \in L(\Delta,\alpha\numdiv^2)$, so the claim is proved.
Actually, the extra factor $\numdiv(n)$ is not very troublesome,
because $\numdiv(n) = o(n^r)$ for every $r>0$
(see~\cite[Section~18.1, Theorem~15, p.~260]{Ha-Wr}).
\end{remark}

%%%%%%%%%%%%%%%%%%%%%%%%%%%%%%%%%%%%%%%%%%%%%%%%%%%%%%%%%%%%%%%

\begin{thebibliography}{9}

\bibitem{Ha-Wr} G. M. Hardy and E. M. Wright,
\textit{An Introduction to the Theory of Numbers}, fifth ed.,
Oxford Univ. Press, 1979.

\bibitem{Ri} T. J. Rivlin,
\textit{Chebyshev Polynomials:
From Aproximation Theory to Algebra and Number Theory}, second ed.,
Wiley, 1990.
%% La primera edici\'on es de 1974, y se titulaba
%% "The Chebyshev Polynomials".

\end{thebibliography}

\end{document}

